\documentclass[12pt]{article}

\usepackage{fontawesome}
\usepackage{hyperref}
\usepackage{xurl}
\usepackage[tagged, highstructure]{accessibility}
\usepackage{bookmark}
\usepackage{enumitem}

\hypersetup{
    colorlinks=false,
    pdfborder={0 0 0},
    pdftitle={Final Project},
    pdfauthor={Adrianna Holden-Gouveia},
    pdfsubject={Introduction to Data - Final Project},
    pdflang={en-US},
    pdfkeywords={data science, presentation, research, final project, Brightspace}
}

% Configure PDF bookmarks for navigation
\bookmarksetup{
    numbered,
    open,
}

% Configure list spacing for better accessibility
\setlist{nosep}


\title{Final Project}
\author{
        Adrianna Holden-Gouveia \\
        Website: \href{https://aholdengouveia.name}{aholdengouveia.name}\\ 
        \date{\vspace{-5ex}}
        %Email: \href{mailto:admin@aholdengouveia.name}{admin@aholdengouveia.name} \\
%        \faLinkedin{: aholdengouveia} \\
        \faGithub {: aholdengouveia} \\
        %\faTwitter {: aholdengouveia} \\
        }

%S\date{\today}


\begin{document}

\maketitle

\section*{Accessibility Notice}
This document is also available in HTML format at:

Link: \href{https://aholdengouveia.name/IntroData/labs/finalproject.html}{\textbf{\nolinkurl{https://aholdengouveia.name/IntroData/labs/finalproject.html}}}

The HTML version provides enhanced accessibility features including keyboard navigation, screen reader support, responsive design, dark mode support, and high contrast options.

\section*{Learning Outcomes}
By the end of this final project, you will be able to:
\begin{itemize}
    \item Research a data topic that wasn't covered in class, using reliable sources checked with the CRAAP test
    \item Plan a lesson for a specific audience, explaining a data topic in a balanced way, with language and examples that fit them
    \item Present a recorded lesson with a clear introduction, supporting points, conclusion, and visuals, within a time limit
    \item Give specific, constructive feedback on other students' presentations
\end{itemize}

\section*{Final Project Instructions}
Pretend you've been asked to teach a short lesson about how data is used. Pick a topic from my Final Project Topics list, linked under Resources, pick who you're teaching, and prepare a 9 to 11 minute lesson. You'll record it, and share it with the class on Brightspace.

The topic has to be new. Don't repeat a lesson from this class, such as collecting data, data integrity, databases and SQL, UI and UX, data visualization, data security, data privacy laws, or big data tools. Your topic and presentation must be about data.

None of these topics are an invitation to be rude to anyone. Cover every topic in a balanced, non-biased way, and focus on the data, not your opinion of the topic. Even politics!

\subsection*{Pick your audience}
Choose who you're teaching, and say who it is at the start of your presentation. For example: other college students, high school students, family members, a sports team, or new employees at a company. Your audience changes your examples, the words you use, and what they need to know.

\subsection*{What your lesson must include}
\begin{itemize}
    \item Why this topic matters to your audience
    \item What data is collected, where it comes from, and how it's used
    \item At least one real example, such as a news story, a company, or an app
    \item What your audience should know or do, such as questions to ask, settings to check, or tools to try
    \item Where they can go to learn more
    \item A clear introduction and conclusion
\end{itemize}

\subsection*{Recording rules}
\begin{itemize}
    \item Your video must be 9 to 11 minutes long.
    \item \textbf{It must be you presenting, in your own voice.} You don't have to show your face. AI voices, AI-generated presenters, and AI-written scripts aren't allowed.
    \item \textbf{Record with a screen capture}, so your slides and anything you show are on screen while you talk. Good free options: OBS Studio, a Zoom meeting where you share your screen and record, or the Screenshot app on a Mac (press Shift, Command, and 5).
    \item \textbf{Don't use PowerPoint's built-in recording}, such as Record Slide Show, or exporting your slides as a narrated video. These won't be accepted.
    \item Don't play other people's videos as part of your presentation. The presentation has to be you.
\end{itemize}

\section*{Steps}
Due dates and the grading rubric are in Brightspace.
\begin{enumerate}
    \item \textbf{Claim your topic.} In the Final Project discussion on Brightspace, start a thread. Make the title your name and your topic. In the body, write who your audience is, a short summary of what you'll cover, and your 5 sources with full links. A topic is only claimed when your thread has all of these. Threads can't be edited, and topics are first come, first served, going by Brightspace's time stamp. If someone already has your topic, pick another one.
    \item \textbf{Record and post your video.} Reply to your own thread with your video, so everyone can watch it and comment.
    \item \textbf{Comment on everyone else's video.} Reply to every other student's presentation. Be kind, and be specific: How well did they research and cite their sources? Was the presentation clear and right for their audience? What's one thing they did well, and one thing that would make it better? Comments like ``good job!'' or ``nice work!'' won't get credit.
    \item \textbf{Turn in your final project} in the Brightspace assignment, with everything listed under Deliverables. Posting your video to the discussion isn't the same as turning in the assignment. You still have to turn it in.
\end{enumerate}

\section*{Deliverables}
Turn in, in the Brightspace assignment:
\begin{itemize}
    \item Your video, or a link to it
    \item Your slides
    \item A report with your audience, your 5 or more sources with citations, and a CRAAP test checklist or paragraph for each source. Presentations based on unreliable sources won't get credit.
    \item Screenshots of every comment you made on other students' presentations
\end{itemize}

\section*{Resources}
My Final Project Topics list
Link: \href{https://www.aholdengouveia.name/IntroData/labs/finalproject-topics.html}{\textbf{\nolinkurl{https://www.aholdengouveia.name/IntroData/labs/finalproject-topics.html}}}

My guide to evaluating sources
Link: \href{https://www.aholdengouveia.name/Howtos/Sourceeval/sourceEvaluation.html}{\textbf{\nolinkurl{https://www.aholdengouveia.name/Howtos/Sourceeval/sourceEvaluation.html}}}

The CRAAP test, with questions to ask about each source
Link: \href{https://web.archive.org/web/20250814000638/https://guides.lib.uchicago.edu/c.php?g=1241077&p=9082343}{\textbf{\nolinkurl{https://web.archive.org/web/20250814000638/https://guides.lib.uchicago.edu/c.php?g=1241077&p=9082343}}}

OBS Studio, free screen recording software for Windows, Mac, and Linux
Link: \href{https://obsproject.com/}{\textbf{\nolinkurl{https://obsproject.com/}}}

Apple: how to record the screen on a Mac
Link: \href{https://support.apple.com/en-us/102618}{\textbf{\nolinkurl{https://support.apple.com/en-us/102618}}}

\end{document}
